\section{Exact replay and release boundary}
\label{sec:replay}

The proofs in Sections~\ref{sec:source-category}--\ref{sec:counterpacket}
are the theorem evidence.  Exact computation is used as a reproducibility
layer for finite arithmetic, matrix traces, and hostile mutations; it does
not replace the all-\(n\) recurrence proof.  No temporary or unpackaged
scout file is a release source or theorem input.

The project-local release chain has four roles:
\begin{enumerate}[leftmargin=*]
\item the producer emits one canonical certificate for the group, descent,
rank, and character ledgers;
\item an independent checker reconstructs the semantic claims rather than
trusting producer booleans;
\item targeted mutations must fail when the support parity, residue
determinant ratio, weight separation, coefficient orbit, or scope firewall is
corrupted;
\item the persistent scoped manifest binds exactly seven release code files
and four release result files while excluding both manifest files; the
44-entry full-project inventory includes it, while the implementation commit
remains a separate provenance stage.
\end{enumerate}

The required exact checks include the following.
\begin{center}
\small
\begin{tabular}{@{}p{0.29\textwidth}p{0.61\textwidth}@{}}
\toprule
Gate & Required result\\
\midrule
Universal group
& Symbolic closure parity, \(6n\) exhaustive count, generator relations,
and independent brute force in the smallest rows\\
Rational form
& \(\delta(r)=r^{-1}\), \(\delta(s)=rs\), exactly two fixed geometric
elements, and distinct Reynolds/transfer denominators\\
Denominator theorem
& Separate-rail survivors exactly \(2,4\); the total-rank mutation must be
rejected for accepting \(n=3\)\\
Third-row character
& 27 monomials, seven independent relations, quotient dimension 20,
residue-corrected group law, both trace vectors, and orbit block \((3,4)\)\\
Scope
& Nonzero class \(\one-\chi_{K/\Q}\) with zero restriction and rank;
all global-root and analytic-promotion flags remain false\\
\bottomrule
\end{tabular}
\end{center}

The tested local promotion protocol is rollback-atomic and exception-safe
under injected failures after its write stages.  This is not a power-loss
atomicity or storage-durability claim.

\subsection{Release-candidate verification record}

The upstream dependency gate reads HCS-C53 from committed git object
\texttt{9d509d3b...} and verifies both the committed Route digest and its
exact artifact tuple.  The unabbreviated commit, digest, and dependency tuple
are recorded in the project source audit; mutable working-tree bytes are not
accepted as provenance.

The immutable project-local release-candidate replay passes 36 of 36 semantic
checks and 93 of 93 unit tests.  The latter include all targeted hostile mutations and an
exhaustive rebound sweep of 198 semantic leaves.  The certificate has 1,078
scalar leaves: 198 semantic, 876 exact-derived, and four allowlisted
history leaves.  The exact artifact hashes are:
\begin{center}
\scriptsize\ttfamily
\begin{tabular}{@{}ll@{}}
payload & f068d5e11ea8e6245e04bd3a30e77140267f835c4e07412ce2009c7fb04ceae1\\
certificate & 780cc9f249e836d3fa5b51a00fd2cdb9af0eac595d929cd1be4d728df1921846\\
independent check & 160b3a9d11354b41404642a3dd22d6e43f2ce576126acb21eb0133e552fc0c0a\\
schema & 4cee6c2252d5743ca3c5fee40ec98fbc945223312d2196fb63a43730281deedf\\
code/results manifest & 62f67e6d4929496974020febab3bc0e2cff45ed153b0cef51937031863d866ba\\
\end{tabular}
\end{center}
These are release-candidate identifiers for the persistent scoped
code/results lane.  The displayed manifest hash is the digest of
\texttt{results/CODE\_RESULTS\_HASHES.sha256}, whose 11 entries exclude both
manifest files; the full-project inventory includes that file.  The
implementation commit is intentionally absent and belongs to a later
provenance stage.  The clean LaTeX build and PDF freeze are performed against
the displayed immutable tuple; later provenance backfill must not alter the
frozen manuscript.

\subsection{Claim boundary}

The unconditional universal result is the projective monomial
equation/group classification.  The packet theorem applies only to
packet-admissible smooth rows, with inherited certified packet data in rows
\(2,3,4\).  The \(n=3\) common-character theorem is first over \(K\).
Nothing in the replay layer upgrades these quantifiers, identifies the
projective monomial stabilizer with the full PGL automorphism group, or
supplies an inert/global root, automorphy, continuation, a functional
equation, or RH.
